%% ma: L11-B15-0006
%% lop: 11
%% chuong: 5
%% bai: 15
%% dang: 11.15.6
%% loai: TL
%% muc: VD
%% dap_an: "-"
%% nguon: "tai-lieu-hoc-tap-mon-toan-lop-11-theo-chuong-trinh-giao-duc-pho-thong-2018.pdf (Lớp Toán Thầy Duy, Phạm Lê Duy, THPT Chu Văn An), Chương 3 Bài 1, Dạng 3, Ví dụ 6 (tr.310)"
%% kien_thuc: "Lớp 11 Bài 15; Bài 7 (tổng n số hạng đầu của cấp số nhân)"
%% trang_thai: chinh-thuc
%% ghi_chu: "Đã duyệt 03/11/2026."
\begin{ex}
Tính các giới hạn sau:

a) $\lim\dfrac{1+3+3^2+\cdots+3^n}{2\cdot3^{n+1}+2^n}$;

b) \[\lim\dfrac{1+\dfrac13+\left(\dfrac13\right)^2+\cdots+\left(\dfrac13\right)^n}{1+\dfrac25+\left(\dfrac25\right)^2+\cdots+\left(\dfrac25\right)^n}.\]
\loigiai{\sloppy a) $1+3+\cdots+3^n=\dfrac{3^{n+1}-1}{2}$. Khi đó biểu thức bằng \[\dfrac{3^{n+1}-1}{2\left(2\cdot3^{n+1}+2^n\right)}=\dfrac{1-\left(\frac13\right)^{n+1}}{2\left(2+\frac23\left(\frac23\right)^{n}\right)}\to\dfrac{1}{4}.\]

b) Tử là tổng $n+1$ số hạng của cấp số nhân công bội $\frac13$: $\dfrac{1-\left(\frac13\right)^{n+1}}{1-\frac13}\to\dfrac32$. Mẫu: $\dfrac{1-\left(\frac25\right)^{n+1}}{1-\frac25}\to\dfrac53$. Vậy giới hạn bằng $\dfrac{3/2}{5/3}=\dfrac9{10}$.}
\end{ex}
